\section{Introduction}
\label{sec:intro}

An agent with persistent memory is an agent that will eventually read two records that disagree. A budget was
revised, a deadline moved, a staging port changed; the store holds both the old statement and the new one, and
sometimes holds two statements that nothing in the record metadata can order. What the agent does at that moment
determines whether persistent memory is an asset or a liability.

The common engineering answer is to label records by reliability --- derived from when each was written --- and to
enforce those labels at the terminal action, so that the agent cannot submit a value the labelling says is stale.
This paper reports that, for small open-weight agents, both halves of that answer can misfire, and measures by how
much.

\subsection{Two failures that motivate the work}

The study began from a specific failure observed during development, not from a general belief that memory
safeguards help. In an exploratory pilot on development splits we gave small agents reliability labels derived only
from record metadata --- the date a record was written. When that metadata could not order the records, as for a
same-day correction or a pair ordered only by dates stated in the text, the labels declared the records
``conflicting''. The agents followed the label instead of the text. For one 7B model on same-day corrections,
correct runs fell from 22 of 24 without labels to 3 of 24 with metadata labels. Enforcing the labels at the terminal
action did not repair this, because enforcement reads the same labels.

A follow-up arbiter that ordered records from textual evidence removed much of that harm on held-out development
phrasings, but exposed the counterpart risk. When the arbiter itself resolved a genuine conflict wrongly --- a
negated cue read as a correction --- enforcement turned five or six of six correct \cond{conflict} answers per model
into arbitrary picks. We call this \textbf{false-resolution harm}: a mechanism confidently asserts an ordering the
evidence does not support, and the agent, or the harness acting for it, commits to a wrong value.

These two failures have the same shape. A reliability signal is produced from the wrong information, and the
machinery downstream --- the agent's deference, or the policy's enforcement --- amplifies it.

\subsection{Approach}

We test \textbf{evidence-basis arbitration}. A deterministic arbiter, Evarb-0.2, assigns every record in a group of
same-attribute records a status together with the \emph{evidence basis} for that status --- an explicit correction
naming the rival value, a stated validity date, the record date, an authoritative source, a genuine contradiction,
or insufficient evidence --- and a \emph{strength}: strong, moderate or none. The agent is shown these labels. At the
terminal action a policy may send one reversible notice when the submitted answer disagrees with the arbitration,
but it \emph{enforces} only resolutions of strong strength. This is \textbf{gating}. If the agent's response to a
notice is unusable, the pre-notice answer stands (\textbf{safe submission}). The full condition, \cond{EBG}, also
asks the agent to search memory before submitting when it has not done so.

The central question has two halves, and we test both:

\begin{quote}
Does evidence-basis-aware arbitration improve persistent-memory reliability in tool-using LLM agents, and can
indiscriminate enforcement of memory labels cause more harmful answers despite improving raw accuracy?
\end{quote}

\subsection{Contributions}

\begin{table}[h]
\centering
\caption{What this paper contributes, and where each claim is supported. ``Confirmatory'' means pre-registered and
frozen before the data existed; everything else is labelled as it is reported.}
\label{tab:contributions}
\small
\begin{tabular}{@{}p{0.46\linewidth}p{0.26\linewidth}p{0.22\linewidth}@{}}
\toprule
\textbf{Contribution} & \textbf{Status of the evidence} & \textbf{Where} \\
\midrule
A pre-registered confirmatory comparison of gated evidence-basis enforcement against enforced record-and-temporal
labels, advisory temporal labels, an LLM arbiter under the same policy, and unlabelled memory, with Holm
correction over the four contrasts.
& Confirmatory & \Cref{sec:results}, \Cref{fig:forest} \\[2pt]
A measurement of the cost of enforcing labels the evidence does not support: harmful-answer and false-resolution
rates per condition, including the pre-registered harm test.
& Confirmatory (harm test); descriptive rates & \Cref{sec:results}, \Cref{fig:tradeoff} \\[2pt]
Component ablations that separate labels, enforcement, gating, safe submission and search-before-submit --- and
show that gating \emph{costs} accuracy.
& Exploratory, BH-adjusted & \Cref{sec:ablations} \\[2pt]
A deterministic-versus-LLM arbiter comparison under one identical downstream policy, separating arbitration
quality from policy.
& Exploratory & \Cref{sec:ablations} \\[2pt]
A sealed evaluation set whose gold is fixed before any text exists, with blind model annotation; released with
hashes. We do not present it as a benchmark of record.
& Artifact & \Cref{sec:dataset} \\[2pt]
A pre-stated account of the mechanism's known failures, written before any confirmatory outcome, extended
afterwards without softening.
& Pre-registered, then extended & \Cref{sec:limitations} \\
\bottomrule
\end{tabular}
\end{table}


The mechanism was frozen before the confirmatory data set was built, and the analysis was pre-registered and frozen
before any confirmatory run. The campaign produced 57{,}888 agent runs with no missing pair, no duplicate and no
harness error.

\textbf{Results.} All four pre-registered contrasts are supported: \cond{EBG} raised task success over enforced
record-and-temporal labels by $17.1$\pp, over advisory temporal labels by $14.5$\pp, over an LLM arbiter under the
same gated policy by $15.9$\pp, and over unlabelled memory by $15.3$\pp, each with a bootstrap interval clear of
zero after Holm correction and in the same direction for all six models. The harmful-answer rate fell from 22.6\% to
3.1\% and arbitrary picks from 252 runs to 13.

\textbf{And two results that cut against the proposed design.} Enforcing every arbitration instead of only
strong-strength ones scored $9.7$\pp{} \emph{higher} while more than doubling harm; and the same gating driven by an
LLM arbiter gave no gain and $8.1$\pp{} more harm. We report these as prominently as the supportive results, because
together they locate the contribution more precisely than the headline contrasts do: what matters is not
enforcement, and not gating, but having an evidence basis that tells enforcement where to stop.

\subsection{Scope}

The study concerns small open-weight models (3B--14B) in one synthetic tool-use harness with templated records,
English text and ISO-format value dates. It makes no claim about frontier models, natural memory logs, implicit
conflicts without lexical cues, or other harnesses. \Cref{sec:limitations} lists the limitations that were fixed
before the confirmatory results existed, and those the results themselves revealed.
